%%%%%%%%%%%%
%
% $Autor: Wings $
% $Datum: 2019-03-05 08:03:15Z $
% $Pfad: SetUp.tex $
% $Version: 4250 $
% !TeX spellcheck = en_GB/de_DE
% !TeX encoding = utf8
% !TeX root = manual 
% !TeX TXS-program:bibliography = txs:///biber
%
%%%%%%%%%%%%

\chapter{Inbetriebnahme}

\section{Standortwahl und Vorbereitung}

\begin{itemize}
	\item \textbf{Untergrund:} Wählen Sie eine ebene, stabile und trockene Arbeitsfläche (Tisch oder Werkbank). Kontrollieren Sie ob das Tischförderband gerade steht. Schiefstellung kann dazu führen, dass die Transportboxen zur Seite rutschen oder herunterfallen.
	
	\item \textbf{Platzbedarf:} Achten Sie darauf, dass an beiden Enden des Transportbands genügend Platz für das Auflegen und Entnehmen der Transportboxen ist.
	
	\item \textbf{Umgebung:} Stellen Sie sicher, dass keine losen Gegenstände oder Kabel in die Nähe der Laufrollen gelangen können. Stellen Sie sicher, dass keine Werkzeuge oder Verpackungsreste auf dem Transportband liegen.

\end{itemize}

%\section{Mechanischer Aufbau}

%\begin{itemize}
%	\item \textbf{Ausklappen:} Klappen Sie die Stützbeine nacheinander vollständig aus.
	
%	\item \textbf{Ausrichtung:} Kontrollieren Sie ob das Tischförderband gerade steht. Schiefstellung kann dazu führen, dass die Transportboxen zur Seite rutschen oder herunterfallen.
	
%	\item \textbf{Mittigkeit:} Prüfen Sie, ob das Transportband mittig auf den Laufrollen liegt.
	
%	\item \textbf{Spannung:} Drücken Sie leicht mit einem Finger von unten mittig an das Transportband. Es sollte stramm sitzen, aber eine minimale Flexibilität von höchstens  0,5$-$1,0 cm haben.
	
%   \ref{sec:Vorspannung einstellen}  Vorspannung anpassen
	
%	\item \textbf{Fremdkörper:} Stellen Sie sicher, dass keine Werkzeuge oder Verpackungsreste auf dem Transportband liegen.
	
%\end{itemize}

\section{Elektrischer Anschluss}

\begin{itemize}
	\item \textbf{Sichtprüfung:} Überprüfen Sie das Gehäuse des Bedienelements und das Netzteil auf sichtbare Schäden.
	
	\item \textbf{Kontrolle Hauptschalter:} Stellen Sie sicher, das der Hauptschalter am Bedienelement in der Postition \fbox{O} steht. 
	
	\item \textbf{Verbindung:} Stecken Sie den Netzteilanschluss des Netzteils in die Buchse am Bedienelement des Tischförderbandes
	
	\item \textbf{Netzanschluss:} Stecken Sie das Netzteil in eine Steckdose (100$-$240 V). Achten Sie darauf, dass das Netzkabel nicht unter Zug steht.
	
\end{itemize}

\section{Inbetriebnahme}
\label{sec:Inbetriebnahme}

% siehe Flowchart Erstes Einschalten

\begin{itemize}
	\item \textbf{Erstes Einschalten: } Nun kann das Tischförderband zum ersten Mal eingeschaltet werden.
	Stellen Sie zum Einschalten den Hauptschalter auf Position \fbox{I}, das System startet.
	Auf der Anzeige ist der Startbildschirm zu sehen. 
	
	\item \textbf{Freigabe der Lichtschranken:} Die Anzeige fragt Sie nach jedem Einschalten, ob beide Lichtschranken freigegeben sind. Stellen Sie erneut sicher, dass keine Transportboxen oder Fremdkörper auf dem Transportband oder zwischen den Lichtschranken liegen.\\
	Bestätigen Sie die Meldung mit dem Quittiertaster. Die Anzeige zeigt nun das Hauptmenü.
	
	Eine übersicht der Anzeige und Menüs finden Sie im Kapitel \ref{sec:Funktionen} \textbf{Funktionen} unter \textbf{Anzeige}.
	
	%\item \textbf{Funktionsprüfung:}
	%
	%Braucht weitere Bearbeitung
	%
	%\begin{itemize}
	%	\item Läuft das \textbf{Transportband} ruhig und ohne schleifende Geräusche?
	%	
	%	\item Bleibt das \textbf{Transportband} in der Mitte oder wandert er zu einer Seite aus? (Falls ja: Justierschrauben an den Laufrollen vorsichtig anpassen).	
	%	
	%	\item \textbf{Modus wählen:} Wählen sie für die Funktionsprüfung mit dem Modusschalter den Modus L aus.  
	%	
	%	\item \textbf{Geschwindigkeit auf Minimum:} Drehen Sie den Drehregler auf die kleinste Stufe.
	%	
	%	\item \textbf{Starten:} Bitte legen sie eine \textbf{Transportbox} auf die rechte Seite des \textbf{Transportbands}, in die \textbf{Lichtschranke 2}. 
	%	
	%\end{itemize}
	%
	%\item \textbf{Geschwindigkeitstest:} Erhöhen Sie langsam die Geschwindigkeit, um die Stabilität der \textbf{Stützbeine} bei Vibrationen zu testen.
\end{itemize}



\section{Betrieb}

\begin{itemize}
	\item \textbf{Aufsichtspflicht:} Bleiben Sie während des Betriebs immer in Reichweite des Hauptschalters.
	
	\item \textbf{Funktionen:} Technische Details und Bedienhinweise finden Sie im Kapitel \ref{sec:Funktionen} \textbf{Funktionen}.
	
	\item \textbf{Not-Aus:} Im Falle einer Blockade stellen Sie sofort den Hauptschalter auf \fbox{O}.
	%Greifen Sie niemals bei laufendem Motor in das Transportband.
	
\end{itemize}
